% Full-page figure: CFL stability diagram for FTCS on the heat
% equation. Two panels (stable and unstable) plus an amplification-
% factor curve. Demonstrates the working signature of the
% explicit-scheme blow-up.
\begin{figure}[p]
\centering

% --- Top panel: stable run at r = 0.4 ---
\begin{tikzpicture}
\begin{axis}[
  width=12cm, height=5cm,
  title={\small (a) FTCS at $r = 0.4$: stable evolution},
  xlabel={position $x$},
  ylabel={$u$},
  xmin=0, xmax=1,
  ymin=-0.1, ymax=1.1,
  grid=both,
  major grid style={engblack!20},
  legend pos=north east,
  legend cell align=left,
  font=\scriptsize,
  samples=120, domain=0:1,
]
% Decaying sin(pi x) at three times, decay rate alpha pi^2 = pi^2
\addplot[engblack, very thick, dashed] {sin(deg(pi*x))};
\addlegendentry{$t = 0$}
\addplot[engblue, very thick]          {sin(deg(pi*x)) * exp(-pi^2 * 0.05)};
\addlegendentry{$t = 0.05$}
\addplot[engred, very thick]           {sin(deg(pi*x)) * exp(-pi^2 * 0.2)};
\addlegendentry{$t = 0.2$}
\end{axis}
\end{tikzpicture}

\vspace{1.5ex}

% --- Middle panel: unstable run at r = 0.6 (schematic) ---
\begin{tikzpicture}
\begin{axis}[
  width=12cm, height=5cm,
  title={\small (b) FTCS at $r = 0.6$: checkerboard blow-up after a few steps},
  xlabel={position $x_j$ (grid index)},
  ylabel={$U_j^n$},
  xmin=0, xmax=1,
  ymin=-12, ymax=12,
  grid=both,
  major grid style={engblack!20},
  legend pos=north east,
  legend cell align=left,
  font=\scriptsize,
]
% Smooth initial sin (rescaled).
\addplot[engblack, very thick, dashed, samples=100, domain=0:1]
  {sin(deg(pi*x))};
\addlegendentry{$n = 0$}
% Grid points: 21 nodes at x = j/20, with checkerboard pattern
% amplified by factor ~3 (=(2r-1)*4 = 0.4^? rough) at n=8 step.
\addplot[engblue, very thick, mark=*, mark size=1pt] coordinates {
  (0.00, 0)  (0.05, 0.9) (0.10, -1.0) (0.15, 1.1) (0.20, -1.2)
  (0.25, 1.4)(0.30,-1.6) (0.35, 1.8) (0.40,-2.1) (0.45, 2.4)
  (0.50,-2.7)(0.55, 2.4) (0.60,-2.1) (0.65, 1.8) (0.70,-1.6)
  (0.75, 1.4)(0.80,-1.2) (0.85, 1.1) (0.90,-1.0) (0.95, 0.9)
  (1.00, 0)
};
\addlegendentry{$n = 8$}
\addplot[engred, very thick, mark=*, mark size=1pt] coordinates {
  (0.00, 0)  (0.05,  6.0) (0.10, -6.5) (0.15, 7.0) (0.20, -7.5)
  (0.25, 8.0)(0.30, -8.5) (0.35, 9.0) (0.40, -9.5) (0.45,10.0)
  (0.50,-10.5)(0.55,10.0) (0.60, -9.5) (0.65, 9.0) (0.70,-8.5)
  (0.75, 8.0)(0.80, -7.5) (0.85, 7.0) (0.90,-6.5) (0.95, 6.0)
  (1.00, 0)
};
\addlegendentry{$n = 16$}
\end{axis}
\end{tikzpicture}

\vspace{1.5ex}

% --- Bottom panel: amplification factor xi(k) for varying r ---
\begin{tikzpicture}
\begin{axis}[
  width=12cm, height=5cm,
  title={\small (c) Amplification factor $\xi(k) = 1 - 4 r \sin^{2}(k \Delta x / 2)$},
  xlabel={normalised wavenumber $k \Delta x / \pi$},
  ylabel={$\xi$},
  xmin=0, xmax=1,
  ymin=-1.6, ymax=1.05,
  grid=both,
  major grid style={engblack!20},
  legend pos=south west,
  legend cell align=left,
  font=\scriptsize,
  samples=200, domain=0:1,
]
\addplot[engblue, very thick]                {1 - 4 * 0.25 * sin(deg(pi*x/2))^2};
\addlegendentry{$r = 0.25$ (well stable)}
\addplot[engred, very thick]                 {1 - 4 * 0.50 * sin(deg(pi*x/2))^2};
\addlegendentry{$r = 0.50$ (CFL bound)}
\addplot[engblack, very thick, dashed]       {1 - 4 * 0.55 * sin(deg(pi*x/2))^2};
\addlegendentry{$r = 0.55$ (unstable)}
\addplot[engblack, thin, dotted, samples=2, domain=0:1] {-1};
\addplot[engblack, thin, dotted, samples=2, domain=0:1] {1};
\end{axis}
\end{tikzpicture}

\caption[CFL stability of FTCS for the heat equation]{The
Courant-Friedrichs-Lewy stability boundary for FTCS applied to the
heat equation $u_t = \alpha u_{xx}$. Panel (a): a run at $r =
\alpha \Delta t / \Delta x^{2} = 0.4$ evolves smoothly, with the
fundamental $\sin(\pi x)$ mode decaying at the correct rate
$\alpha \pi^{2}$. Panel (b): a run at $r = 0.6$ develops a
short-wavelength checkerboard pattern within the first dozen steps
and within roughly twenty steps the divergence overwhelms the
solution by orders of magnitude. Panel (c): the amplification
factor $\xi(k) = 1 - 4 r \sin^{2}(k \Delta x / 2)$ shows the
mechanism. At $r = 0.5$ the curve just touches $\xi = -1$ at the
Nyquist wavenumber $k \Delta x = \pi$, the marginal case. At $r =
0.55$ the curve drops below $-1$ at the Nyquist, so $|\xi| > 1$ for
the highest-wavenumber mode; that mode amplifies by $|\xi|^{n}$ per
$n$ steps, producing the checkerboard. The CFL condition
$r \le 1/2$ is the algebraic statement of $|\xi(k)| \le 1$ across
the full Nyquist range. \cite{paper:v2ch12-courant-friedrichs-lewy-1928}}
\label{fig:vol02:ch12:cfl-stability}
\end{figure}
